% Instructor-authored original practice bank. The LSAT and LSAC names are used
% only to identify the source of skill inspiration. This material is not copied
% from, affiliated with, approved by, or endorsed by LSAC.

\renewcommand{\labelenumii}{(\Alph{enumii})}

\section*{Reading Comprehension: Original Practice Bank}

\begin{center}
\fbox{\parbox{0.92\linewidth}{\textbf{Original-material notice.} Every passage,
fact pattern, question, and answer choice below is instructor-authored original
material. The names LSAT and LSAC are used only to identify skill inspiration.
This practice bank is not an official test form and is not endorsed by LSAC.}}
\end{center}

Each packet is designed for approximately 30--50 minutes of reading, individual
reasoning, and group discussion. The source-map comments identify only the
disclosed administration, location, and broadly tested reading skill; they do not
reproduce the source material.

\subsection*{Packet 1: Listening for Ecological Change (Questions 1--5)}

For decades, city ecologists counted birds mainly by sending trained observers
along fixed routes at dawn. That method remains valuable, but it limits both the
number of sites and the times that can be sampled. A newer approach places small
audio recorders on rooftops and in parks. Software proposes labels for the calls
in the recordings, and human reviewers confirm uncertain cases. Because the
recorders can operate every night and morning, they can reveal changes that a
weekly observer might miss.

The method is sometimes described as automated monitoring, a description that
can obscure how much judgment remains involved. A microphone records machinery,
wind, and overlapping calls as readily as it records a clear song. Software
trained on rural recordings may misclassify calls distorted by traffic noise.
Even a correct species label does not establish how many birds were present: one
loud individual can call repeatedly, while a flock may pass silently. Thus, the
recordings do not replace field observation. They extend it.

This distinction matters when researchers use sound archives to reconstruct
change. Suppose a species appears more often in recent recordings than in older
ones. The difference might reflect a growing population, but it might instead
reflect improved microphones or a recorder moved closer to vegetation. A strong
study therefore preserves equipment notes, placement maps, weather records, and
samples that humans can audit. With those controls, an archive can do something
a traditional survey rarely can: allow a later researcher to ask a question that
the original collector never anticipated.

The greatest contribution of acoustic monitoring may consequently be neither
speed nor automation, but revisability. A written checklist records an observer's
conclusion. An audio file also preserves some of the evidence from which the
conclusion was drawn. Future listeners can challenge a label, search for a newly
recognized call, or compare changes in the surrounding soundscape. The archive is
not neutral---what was recorded still depends on microphone placement and
institutional choices---but it keeps more of the inferential path open to review.

\begin{enumerate}
% Source map: January 2023 | Reading Comprehension | Q1 | weaken an application
\item A research team claims that a rise in recorded calls proves that a bird's
population increased. Which fact most weakens that claim under the passage's
analysis?
  \begin{enumerate}
  \item The newer recordings were made with more sensitive microphones placed
        closer to nesting vegetation.
  \item The species has more than one type of call.
  \item The recordings were stored in a common digital format.
  \item Human reviewers agreed about most proposed labels.
  \item Dawn surveys are usually conducted before breakfast.
  \end{enumerate}

% Source map: January 2023 | Reading Comprehension | Q2 | main point
\item Which statement best expresses the passage's main point?
  \begin{enumerate}
  \item Acoustic archives should replace all direct bird surveys.
  \item Traffic noise makes urban ecological research impossible.
  \item Audio monitoring is most useful when treated as revisable evidence that
        complements observation and is documented carefully.
  \item Software identifies urban bird calls more accurately than experts do.
  \item Population size can always be inferred from the number of calls recorded.
  \end{enumerate}

% Source map: January 2023 | Reading Comprehension | Q3 | supported inference
\item The passage most strongly supports which inference?
  \begin{enumerate}
  \item A later research question may be investigated with an old audio archive
        even if the original study did not ask that question.
  \item Written checklists contain no useful ecological information.
  \item Rural recordings are always clearer than urban recordings.
  \item A silent flock cannot be observed by any method.
  \item Equipment notes matter only when software labels are wrong.
  \end{enumerate}

% Source map: January 2023 | Reading Comprehension | Q4 | function of a detail
\item The example of one loud bird calling repeatedly primarily serves to
  \begin{enumerate}
  \item show why cities should prohibit rooftop microphones.
  \item illustrate why call frequency cannot automatically be equated with the
        number of birds present.
  \item establish that solitary birds are louder than flocks.
  \item explain why weather records are unnecessary.
  \item prove that human reviewers cannot identify birds.
  \end{enumerate}

% Source map: January 2023 | Reading Comprehension | Q5 | passage detail
\item According to the passage, which practice helps researchers distinguish an
ecological change from a change in recording conditions?
  \begin{enumerate}
  \item Discarding recordings after labels are assigned
  \item Counting every detected call as a separate bird
  \item Preserving equipment and placement records
  \item Using only one morning of data per year
  \item Preventing later researchers from revising labels
  \end{enumerate}
\end{enumerate}

\subsection*{Packet 2: What a Route Map Should Show (Questions 6--13)}

\textbf{Passage A.} A transit map is an argument about relevance. Riders usually
need to know the order of stops, where lines connect, and which services require
a transfer; they rarely need a scale drawing of every curve in the track. A
schematic map therefore distorts physical distance to make network structure
legible. Critics sometimes call such distortion inaccurate. Yet a map that
preserves geography perfectly can be practically misleading if central stations
are crowded together until their labels cannot be read. Accuracy must be judged
relative to the task.

Digital trip planners appear to escape this choice by calculating a route for
each user. But their apparent precision introduces a different abstraction. A
planner ranks options using assumptions about walking speed, transfer risk, and
the cost of waiting. If those assumptions remain invisible, a route displayed as
``best'' may be inexplicable. Designers should therefore reveal consequential
assumptions and let riders adjust them. A useful map simplifies; a trustworthy
one also shows users how the simplification matters.

\textbf{Passage B.} Personalized routing can improve access, but customization is
not a substitute for a stable public diagram. A shared map gives riders a common
language: ``two stops beyond the river'' or ``change at the square'' works even
when phones fail or travelers use different applications. It also permits people
to form a mental picture of the system rather than follow isolated instructions.

For that reason, agencies should resist displaying only individually generated
routes. They should maintain a common schematic and test it with riders who have
different visual, linguistic, and mobility needs. The diagram need not preserve
geographic scale. Indeed, enlarging a dense center may improve it. But departures
from geography should be consistent, and accessibility information should not be
hidden behind optional menus. A public map succeeds not by containing every fact,
but by making a defensible set of facts available to everyone.

\begin{enumerate}
\setcounter{enumi}{5}
% Source map: January 2023 | Reading Comprehension | Q6 | shared concern
\item Both passages are primarily concerned with
  \begin{enumerate}
  \item proving that geographic scale must never be altered.
  \item how simplification and design choices affect the usefulness of transit
        information.
  \item replacing transit diagrams with printed timetables.
  \item showing that riders always choose the fastest route.
  \item explaining how rail lines are engineered.
  \end{enumerate}

% Source map: January 2023 | Reading Comprehension | Q7 | inference about Passage A
\item The author of Passage A would most likely agree that
  \begin{enumerate}
  \item a route-ranking system should disclose assumptions that materially affect
        its recommendations.
  \item any distortion of physical distance makes a map inaccurate.
  \item a digital planner never simplifies information.
  \item all riders assign the same cost to waiting.
  \item a map is useful only if it contains every available fact.
  \end{enumerate}

% Source map: January 2023 | Reading Comprehension | Q8 | function of phrase
\item In Passage B, the examples ``two stops beyond the river'' and ``change at
the square'' are used to illustrate
  \begin{enumerate}
  \item a defect caused by consistent diagrams.
  \item the engineering vocabulary of transit employees.
  \item how a shared map supports communication among riders.
  \item why mobile phones never work underground.
  \item a demand that maps preserve geographic scale.
  \end{enumerate}

% Source map: January 2023 | Reading Comprehension | Q9 | underlying principle
\item Which principle is most consistent with both passages?
  \begin{enumerate}
  \item Information design should be evaluated in light of users' tasks and the
        consequences of what is emphasized or omitted.
  \item A representation is reliable only if it reproduces physical space exactly.
  \item Personalized systems should always replace common public resources.
  \item More information necessarily produces a clearer design.
  \item Transit agencies should avoid consulting riders.
  \end{enumerate}

% Source map: January 2023 | Reading Comprehension | Q10 | relationship of purposes
\item Which statement best describes the relationship between the passages?
  \begin{enumerate}
  \item Passage A condemns schematic maps, while Passage B defends them.
  \item Passage A analyzes task-relative accuracy and algorithmic assumptions;
        Passage B emphasizes the continuing value of a shared public diagram.
  \item Passage A discusses only buses, while Passage B discusses only trains.
  \item Passage B supplies experimental data that refute Passage A.
  \item Passage B argues that no transit information should be personalized.
  \end{enumerate}

% Source map: January 2023 | Reading Comprehension | Q11 | view in A but not B
\item Which claim is made in Passage A but not in Passage B?
  \begin{enumerate}
  \item A schematic may alter geographic distances.
  \item A dense city center can require visual enlargement.
  \item Route planners rank choices using assumptions about matters such as
        walking speed and waiting.
  \item Transit information should serve riders' needs.
  \item Designers necessarily intend to deceive users.
  \end{enumerate}

% Source map: January 2023 | Reading Comprehension | Q12 | relationship assessment
\item Passage B's position on shared diagrams most directly qualifies which idea
from Passage A?
  \begin{enumerate}
  \item That track curves are physically real
  \item That digital tools can calculate a route for each individual user
  \item That crowded labels can be difficult to read
  \item That users care about transfers
  \item That assumptions can be consequential
  \end{enumerate}

% Source map: January 2023 | Reading Comprehension | Q13 | comparative attitude
\item The authors' attitudes toward schematic distortion are most accurately
described as
  \begin{enumerate}
  \item uniformly hostile in both passages.
  \item approving in Passage A but hostile in Passage B.
  \item accepting in both passages when distortion improves legibility and is
        responsibly designed.
  \item indifferent in Passage A and unaware in Passage B.
  \item accepting only when the map becomes geographically exact.
  \end{enumerate}
\end{enumerate}

\subsection*{Packet 3: Repairing Public Objects (Questions 14--21)}

Museums once tended to describe restoration as the recovery of an object's
original state. That phrase suggested a single correct destination. In practice,
however, an old object may embody several historically important states. A town
clock, for example, may have acquired a new face after a civic reform and a new
mechanism after a factory fire. Removing every later alteration could erase the
very history that made the clock significant.

Contemporary conservators often begin with a different question: which values
should this intervention preserve? Material stability is one value, but so are
evidence of use, the ability of a community to recognize an object, and the
possibility that future specialists will learn from it. These aims can conflict.
Cleaning may make an inscription legible while removing residue that could later
be analyzed. Replacing a worn handle may allow a machine to operate while hiding
how generations of workers held it.

The answer is not to refuse every intervention. Neglect is itself a choice, and
decay can destroy options just as surely as an aggressive repair. Instead,
conservators increasingly favor documented, limited, and when possible reversible
steps. A reversible adhesive, for instance, lets a future conservator undo a
repair if better evidence emerges. Detailed photographs and treatment notes do
not make an intervention neutral, but they expose its reasoning to scrutiny.

Community consultation complicates the process productively. The people for whom
an object has ongoing ceremonial or practical importance may value continued use
more than a museum values an untouched surface. Consultation does not mean that
every request can be granted: operation may destroy a fragile object, and groups
may disagree among themselves. It does mean that technical expertise alone
cannot identify all relevant values. Restoration is best understood not as a
return journey to a fixed origin, but as a responsible choice among possible
futures, made with evidence, restraint, and accountable judgment.

\begin{enumerate}
\setcounter{enumi}{13}
% Source map: January 2023 | Reading Comprehension | Q14 | main point
\item The main point of the passage is that restoration should
  \begin{enumerate}
  \item always remove later alterations.
  \item be treated as an accountable choice among competing values and possible
        futures rather than recovery of one self-evident original state.
  \item give material stability no weight.
  \item be controlled exclusively by community vote.
  \item avoid all documentation until treatment is complete.
  \end{enumerate}

% Source map: January 2023 | Reading Comprehension | Q15 | detail/inference
\item Why, according to the passage, might removing a later alteration be a loss?
  \begin{enumerate}
  \item Later alterations are always more attractive than original features.
  \item The alteration may itself record an important part of the object's history.
  \item Museums are legally prohibited from revealing original surfaces.
  \item Later materials never decay.
  \item Every community prefers recently altered objects.
  \end{enumerate}

% Source map: January 2023 | Reading Comprehension | Q16 | function of reference
\item The discussion of a reversible adhesive primarily illustrates
  \begin{enumerate}
  \item a repair method that preserves the possibility of later reconsideration.
  \item why all adhesives damage old objects.
  \item a way to eliminate the need for treatment records.
  \item proof that conservation decisions are neutral.
  \item why objects should remain in continuous operation.
  \end{enumerate}

% Source map: January 2023 | Reading Comprehension | Q17 | author agreement
\item The author would most likely agree with which statement?
  \begin{enumerate}
  \item Doing nothing to a deteriorating object is free of consequences.
  \item Technical experts can identify every value relevant to an object without
        consultation.
  \item A limited intervention can be justified even though it is not neutral.
  \item Continued use must always take priority over material survival.
  \item Documentation converts a contested choice into an objective fact.
  \end{enumerate}

% Source map: January 2023 | Reading Comprehension | Q18 | inference about proponents
\item Advocates of the approach favored in the passage would most likely support
  \begin{enumerate}
  \item replacing every worn component with a visually identical new one.
  \item recording the reasons for choosing one treatment over another.
  \item prohibiting future specialists from revisiting a repair.
  \item defining significance solely by market value.
  \item concealing uncertainty from affected communities.
  \end{enumerate}

% Source map: January 2023 | Reading Comprehension | Q19 | passage detail
\item According to the passage, cleaning an inscription can create a conflict
because it may
  \begin{enumerate}
  \item improve legibility while removing material that could yield future evidence.
  \item prevent anyone from reading the inscription.
  \item make community consultation impossible.
  \item guarantee that the object can operate.
  \item reveal that the object has no history.
  \end{enumerate}

% Source map: January 2023 | Reading Comprehension | Q20 | supported inference
\item Which inference is best supported by the passage?
  \begin{enumerate}
  \item Reversibility and documentation can reduce, but do not erase, the stakes
        of intervention.
  \item All groups consulted about an object will assign it the same value.
  \item An object's earliest state is never historically important.
  \item A functioning object cannot also be fragile.
  \item Conservators have stopped caring about material stability.
  \end{enumerate}

\par\medskip

% Source map: January 2023 | Reading Comprehension | Q21 | invoked principle
\item Which principle best captures the author's view of consultation?
  \begin{enumerate}
  \item Expertise should exclude all knowledge derived from use or tradition.
  \item Affected communities contribute relevant knowledge, although their views
        do not mechanically settle every conservation decision.
  \item Consultation is useful only when every participant agrees.
  \item Technical risks should never limit ceremonial use.
  \item Museums should transfer every object requested by any group.
  \end{enumerate}
\end{enumerate}

\subsection*{Packet 4: Designing Instruments by Constraint (Questions 22--27)}

The composer Amira Vale builds pieces around instruments that do not yet exist.
Her collaborators construct modules---reeds, resonating boxes, keys, and digital
sensors---that performers can rearrange. Reviewers often describe the instruments
as devices for producing an unlimited range of sound. Vale rejects that account.
She begins each piece by imposing a narrow mechanical rule: perhaps air may enter
through only one chamber, or every key must also alter the instrument's balance.

These constraints do more than give the instrument a distinctive tone. They make
performance visibly consequential. On a conventional keyboard, pressing one key
usually leaves the layout unchanged. On one Vale instrument, pressing a key
shifts a weight, gradually tilting the resonating box. The performer must decide
not only which pitch is wanted now, but which pitches will remain reachable later.
An audience can often see the accumulating result of those decisions.

Vale's scores are similarly unusual. They specify goals and prohibitions rather
than a complete sequence of notes. One page instructs a trio to reach a stable
chord without allowing any player to repeat a motion. Different performances
therefore sound unlike one another, yet they are not arbitrary: each is a solution
to the same designed problem. The score functions less like a script than like a
set of rules for a game whose moves carry audible costs.

It would be a mistake, however, to identify constraint with deprivation. A useful
constraint creates relationships that a performer can investigate. If a rule
merely removes choices without making the remaining choices interact, it produces
thin music. Vale's achievement lies in designing limits that generate structure.
The instruments are interesting not because they can do anything, but because
each carefully chosen action changes what can happen next.

\begin{enumerate}
\setcounter{enumi}{21}
% Source map: January 2023 | Reading Comprehension | Q22 | main point
\item Which statement best expresses the passage's main point?
  \begin{enumerate}
  \item Vale's work uses generative constraints so that performance becomes a
        sequence of consequential choices rather than unlimited improvisation.
  \item Vale seeks to reproduce conventional keyboards as accurately as possible.
  \item A successful musical instrument must be capable of every sound.
  \item Vale's scores determine every note before a performance begins.
  \item Visible movement prevents an audience from hearing musical structure.
  \end{enumerate}

% Source map: January 2023 | Reading Comprehension | Q23 | conforming scenario
\item Which scenario most closely conforms to Vale's approach as described?
  \begin{enumerate}
  \item A performer may play any sound at any time, with no effect on later options.
  \item A flute has several decorative colors but produces the same fixed sequence.
  \item A percussion frame rotates after each strike, changing which surfaces the
        performer can reach during the next phrase.
  \item A recording plays identically whenever a button is pressed.
  \item A score forbids half the notes but creates no relation among those remaining.
  \end{enumerate}

% Source map: January 2023 | Reading Comprehension | Q24 | analogy
\item The passage's description of Vale's scores is most analogous to
  \begin{enumerate}
  \item a photograph that records a completed event.
  \item a recipe whose quantities and timing may never vary.
  \item a game defined by shared rules but allowing many different sequences of play.
  \item a list of unrelated sounds selected by chance.
  \item a machine that conceals every consequence of a user's action.
  \end{enumerate}

% Source map: January 2023 | Reading Comprehension | Q25 | supported inference
\item The passage most strongly suggests that Vale would criticize a constraint if it
  \begin{enumerate}
  \item made present choices affect later possibilities.
  \item permitted different performances to sound different.
  \item simply reduced the available choices without creating meaningful interaction.
  \item could be perceived by an audience.
  \item required collaboration among performers.
  \end{enumerate}

% Source map: January 2023 | Reading Comprehension | Q26 | primary purpose
\item The primary purpose of the passage is to
  \begin{enumerate}
  \item explain how Vale's instruments and scores use carefully designed limits to
        create musical structure.
  \item argue that all conventional instruments should be destroyed.
  \item present measurements comparing the volume of several instruments.
  \item prove that improvisation requires no rules.
  \item describe the commercial manufacture of digital sensors.
  \end{enumerate}

% Source map: January 2023 | Reading Comprehension | Q27 | most strongly supported
\item Which statement is most strongly supported by the passage?
  \begin{enumerate}
  \item Every performance of a Vale score contains the same sequence of pitches.
  \item Vale regards unpredictability alone as sufficient for musical value.
  \item The modular instruments conceal the effects of performers' decisions.
  \item Two performances can differ substantially while satisfying the same score.
  \item Vale's mechanical rules apply only after a performance ends.
  \end{enumerate}
\end{enumerate}

\section*{Answer Key and Concise Explanations}

\begin{description}
\item[1. A] Better equipment and placement provide an alternative explanation for more detected calls.
\item[2. C] The passage presents archives as documented, revisable evidence that extends field observation.
\item[3. A] Preserved recordings can be reexamined for questions not anticipated initially.
\item[4. B] One bird can generate many calls, so detections do not directly count individuals.
\item[5. C] Equipment and placement records allow changes in conditions to be assessed.
\item[6. B] Both authors analyze how design choices shape usable transit information.
\item[7. A] Passage A explicitly favors revealing consequential ranking assumptions.
\item[8. C] The phrases show how a common diagram supplies shared reference points.
\item[9. A] Both evaluate selective representation by its purpose and consequences.
\item[10. B] The passages emphasize different, compatible aspects of responsible transit design.
\item[11. C] Only Passage A discusses assumptions embedded in route-ranking systems.
\item[12. B] Passage B warns that personalized routes should not displace a shared diagram.
\item[13. C] Both accept purposeful, consistently designed departures from geographic scale.
\item[14. B] The author replaces the single-original-state model with accountable value choices.
\item[15. B] A later change can itself be historically significant evidence.
\item[16. A] Reversibility keeps open the option of revising a treatment later.
\item[17. C] The passage permits limited intervention while denying that it is neutral.
\item[18. B] Documented reasoning is central to the accountable approach defended.
\item[19. A] Cleaning can trade present legibility against material for future analysis.
\item[20. A] These practices preserve options and scrutiny without eliminating judgment.
\item[21. B] Consultation adds relevant values but does not automatically dictate treatment.
\item[22. A] Vale's limits make each choice reshape the field of later choices.
\item[23. C] Each action changes the physical possibilities available next.
\item[24. C] The score supplies common constraints while permitting varied valid performances.
\item[25. C] The author distinguishes generative constraints from mere deprivation.
\item[26. A] The passage explains the structural role of designed limits across instruments and scores.
\item[27. D] Shared constraints can yield performances that sound unlike one another.
\end{description}
